In this final chapter, I step back from the specifics of each study to reflect more broadly on what the thesis contributes. The specific contribution from each separate study is outlined in the corresponding article \textbf{(See Appendix)}. The sections that follow are organized into four sections describing theoretical contributions, managerial implications, methodological notes, and limitations (with directions for future research). Each section connects insights across the empirical studies to identify recurring patterns and broader lessons that extend beyond the individual analyses. Across these sections, the chapter provides an integrated synthesis of the thesis’s main contributions and outlines directions for future research.

\section{Theoretical Contributions}

This thesis contributes to marketing and retail theory by demonstrating how large-scale behavioral data, when analyzed through established consumer behavior theories, can be used to examine how strategic decisions relate to customer responses across different retail settings. Following a consumer-based strategy perspective, which emphasizes the use of insights from customer behavior to inform strategic decisions ((\cite{Hamilton2016-ca}), the studies treat observable customer responses as the unit of analysis. These responses provide a behavioral basis for assessing how strategic decisions play out across contexts, rather than interpreting outcomes in isolation.

First, this thesis highlights the importance of retail context when evaluating the strategies’ effectiveness. Across Studies 1–4, results illustrate that customer responses to retail strategies are not consistent across formats, brands, channels, and platform participants. Studies 1 and 2 show that, even within the same retailer, the impact of pricing strategies and discounts for PL products cannot be generalized. Price elasticities and profit outcomes differ systematically by store format, PL tier, and product category. These findings align with prior research on multiformat retailing and PLs (e.g., \cite{Geyskens2010-bx,Haans2011-cv}), and with broader evidence on variations in marketing mix across formats (\cite{Wichmann2021-vm}). They also suggest that performance should not be evaluated using aggregated effects alone, but must account for differences across context, brand, and assortment. Study 4 extends this context principle to digital platforms to reveal that Spotify’s playlist follower dynamics during the pandemic did not follow uniform patterns but differed by curator type and content composition. It further shows that, consistent with broader platform dynamics, external shocks can shift visibility and outcomes unevenly across platform participants (\cite{Wichmann2022-gq}).

Second, the studies jointly demonstrate that customer responses help explain how strategic actions translate into observable outcomes. Rather than presenting outcomes without explanation, each study links strategic actions (e.g., pricing, channel introduction, and platform management) to customer responses in purchasing or platform interaction. For example, Study 3 draws on shopping utility maximization to explain why online channel introduction for a grocery context altered the mix and monetary value of purchases over time. The study expands on work by \textcite{Campo2015-pt} and \textcite{Melis2016-mf}, and it aligns with recent evidence on household adaptation to economic pressures and category-level adjustments (\cite{Scholdra2022-rp}). It shows that long-term revenue changes are driven by persistent shifts in the categories from which customers purchase. In particular, customers increase spending in categories with higher online delivery convenience, and this leads to higher value per visit over time. Study 4, in turn, shows that customer behavior on digital platforms reflects broader shifts in daily routines during COVID-19. As time spent at home increased, users adjusted their listening behavior, leading to changes in how attention was distributed across playlists. These patterns are more clearly understood when observed behavior is interpreted in relation to external conditions.

Third, the thesis refines existing knowledge on pricing and channel strategies. Through comparison of regular price changes, temporary discounts, and PL tiers, in Studies 1 and 2 the findings show that neither price nor discount elasticity is consistent across products and store formats. For example, Study 2 shows that discount responsiveness is stronger in hypermarkets than in smaller formats. Study 1 further shows that profitability differs across PL tiers depending on store format. These results build on previous findings (e.g., \cite{Ailawadi2006-he,Shankar1996-kq,Widdecke2022-xo}) and highlight that the effectiveness of pricing strategies depends on both store format and assortment structure. Study 3 similarly shows that the effects of online channel adoption on long-term value are not uniform but dependent on how customers adjust their purchasing behavior over time.

Last, the thesis emphasizes both the value and the limits of data-driven approaches. By leveraging unique retailer and platform datasets that directly connect strategy to observed behavior, the studies show that insights do not come from data volume alone but from how data are interpreted. Common generalizations, such as the belief that discounting always increases sales, or that online channels always add value, are shown to depend on context. The thesis therefore substantiates calls for theory-driven, context-aware analytics (\cite{Dekimpe2020-bh}), echoing similar arguments for more adaptive, context-specific decisions on marketing mix (\cite{Wichmann2021-vm}). It suggests a model for future research that balances prediction with explanation.

\section{Managerial Implications}

Many of the findings in this thesis are not new to practitioners. Price and discounting are not the same thing. Store formats influence how customers behave. Platforms tend to favor their own content. However, the issue is that these insights are not often reflected in how decisions are actually implemented or evaluated in practice. As noted by \textcite{McKinsey-Company2021-wc}, performance reports often aggregate prices and promotions under a single pricing which obscures their distinct effects. Store formats have traditionally been treated primarily as alternative offering channels, not as distinct behavioral contexts. Recent industry reports (e.g., \cite{Accenture2024-vg,Bain-Company2025-hz}) suggest that retailers are beginning to rethink this view.

The contribution of this thesis is therefore not to introduce new ideas but to quantify and structurally outline what is already known. By testing these differences across formats, categories, and time, the studies show where common assumptions hold and where they do not. In particular, they show that applying the same strategy across contexts can lead to very different outcomes. This is consistent with prior work showing that managerial insights are more useful when grounded in real-world data and linked to practical decisions (\cite{Schauerte2023-dn}).

Take, for instance, the idea that deep discounts drive sales. This is generally true, but the effect could depend on format, product category, and timing. Studies 1 and 2 show that profit gains from price increases on budget PLs hold in hypermarkets but not in supermarkets. Promotion fatigue also varies across categories, and customers’ past exposure to discounts reduces discounts’ future effectiveness. These findings are not new, but they are rarely quantified or compared across formats, categories, and customer groups. The key implication is not simply that discounts work or that formats matter, but that applying the same pricing approach across products, formats, or customer groups can lead to suboptimal outcomes. Managers therefore need to be more precise in how they frame and evaluate pricing decisions.

In place of such broad advice as “be agile” or “embrace personalization” (advice often found in industry reports; e.g., \cite{Accenture2024-vg,McKinsey-Company2019-cx}), by quantifying effects across formats, categories, and time this thesis shows where differences actually matter. The findings are based on observed customer behavior in real empirical settings, and they demonstrate that practical improvements come from making more precise and structured use of existing data.

This thesis also challenges the common assumption that more data will necessarily improve decision-making. Retailers and platforms already have access to large amounts of data, yet the main challenge often lies in how that data is interpreted. Across the studies, better insights are not driven by larger datasets but by framing decisions through consumer behavior theory (e.g., in terms of utility, perceived value, or platform structure). In this sense, theory can serve as a practical tool for interpreting observed patterns, especially when data are noisy or incomplete. Understanding customer behavior—grasping why customers switch channels, respond to price changes, or interact with content—may be more transferable than relying solely on additional data or tools. This does not imply that data are unimportant, but it does imply that more data alone may not always lead to better decisions.

Further, this thesis contributes is its attention to time. In practice, for example in campaign planning, time is often treated as an operational detail rather than as part of how decisions are made. Across the four studies, the results show that the duration of effects, their persistence over time, is just as important as the effects themselves. Study 1 distinguishes between short- and long-run price responses. Study 3 tracks how customer behavior changes after online channel adoption and which of these changes persist over time. Study 4 examines how follower patterns shift in response to a system-wide disruption. These analyses focus on more concrete questions: what persists, for whom, and for how long? By treating time as part of the analysis, not only as a reporting interval, the thesis highlights how effects can build, fade, or change over time.

Finally, while the studies are empirically grounded in grocery retail and music streaming, their implications extend beyond these settings. The findings from multiformat pricing and multichannel allocation are relevant for other frequently purchased categories such as health, household goods, or personal care. Similarly, the results on platform behavior apply not only to Spotify but also to retail environments that manage digital visibility through search results, home page placements, and recommendation systems. What matters is not the domain but how format, price, assortment, and customer decision-making interact. These interactions are not industry-specific; they arise whenever businesses manage what customers see and how they choose across formats and channels.

\section{Methodological Notes}

Throughout this thesis, theory serves as a guide in shaping the empirical approach. The choice of variables, model structures, and empirical specifications are all grounded in prior research to ensure that the results speak to meaningful constructs in retail strategy and consumer behavior. In particular, each study connects theoretical findings to measurable and observable outcomes (e.g., how price changes affect sales and profit, how channel adoption changes purchasing behavior, or how platform exposure relates to follower dynamics). This helps ensure that the empirical results remain connected to meaningful managerial questions.

A recurring methodological consideration across the studies is the use of panel data structures that support analysis of customer behaviors’ change over time. This thesis uses three complementary empirical strategies depending on the level of decision: error-correction modeling (ECM) for pricing responses, DiD for channel adoption, and panel fixed effects models for platform behavior. These approaches are used not only to control for unobserved differences but also to capture how effects evolve over time. Panel data enables examination of not just whether a strategic action has an effect but how that effect may change over time. Since the research questions emphasize both immediate and longer-term responses to retail decisions, temporal dynamics are treated in the course of the analysis. In Studies 1 and 2, ECM separates short-term effects from longer-term adjustments in sales and profits (\cite{Pauwels2007-ky}). Study 3, comparing pre- and post-periods between matched customer groups, uses a difference-in-difference design to track sustained behavioral changes after digital channel adoption (\cite{Golder2022-bo}). Study 4 applies fixed effects panel models with time controls to isolate change while accounting for shocks related to the COVID-19 pandemic. Across all four studies, these methods allow the analyses to identify whether effects are short-lived or persistent, and how they develop over time. This is consistent with prior research emphasizing that consumer behavior is dynamic and should be analyzed over time (\cite{Zhang2021-os}).

Another important aspect of the methodological approach is the use of real-world data from retailers and digital platforms. This thesis prioritizes high external validity by studying actual observed customer behavior in real decision environments. Across all studies, efforts were made to account for unobserved heterogeneity, to control for confounding factors, and to use robust estimation techniques. The first study applies ECM to estimate both short- and long-term price effects across PL tiers and formats; it thereby links sales responses to profit outcomes. The second study uses a sequential regression design to isolate differential effects of regular pricing and discounts while accounting for brand- and category-level moderators. The third study employs a difference-in-differences design with PSM to estimate long-term behavioral changes following online channel adoption. The fourth study, leveraging government restriction indices and mobility data to contextualize pandemic effects uses a unit fixed effects panel model to assess changes in playlist followers over time. These designs reflect the shared goal of extracting credible insights from observational data while acknowledging its boundaries. This design aligns with recent calls for an empirics-first approach, where theory generation is grounded in the systematic analysis of real-world behavioral data, not imposed a priori (\cite{Golder2022-bo}).

What ties the empirical studies together is the effort to link observed behavioral outcomes such as purchasing responses, channel adoption, and follower changes to broader retailer-level outcomes. This approach reflects the belief that meaningful analysis is derived not from statistical complexity alone but from efforts to build bridges between components of a problem (\cite{Leeflang2014-fs}). For example, in the pricing studies, the analysis moves beyond elasticity to consider profit response and substitution across PL tiers. In the online channel adoption study, the analysis goes from product mix shifts to sustained revenue differences. In the platform study, follower changes are used as observable indicators of exposure and response. This ensures that the analysis remains relevant for understanding performance outcomes.

Finally, the studies follow a consistent structure from theory to empirical specification and interpretation. Each study starts from a theoretical framework to guide the choice of variables and model specification. The empirical results are then interpreted in light of these theoretical expectations. This structure helps maintain comparability across contexts, and it ensures that the analysis remains grounded in both theory and observed behavior. Using established methods with real-world data strengthens the validity of the results and supports their practical relevance.

\section{Limitations and Future Directions}

Like all empirical research, this thesis is shaped by constraints in data, scope, and method. Recognizing these limitations is important, not only for the validity of the current findings but also for guiding future research that builds on and improves them. While each empirical study has its own specific boundaries, several overarching themes emerge across the thesis. These include contextual scope, modeling constraints, and data design and inference.

All four studies are embedded in specific retail and platform environments. This fact strengthens internal validity but limits broader generalization. The pricing analyses (Studies 1 and 2) rely on data from a single Nordic grocery retailer where loyalty norms, assortment structures, and competitive intensity may differ from what is found in other markets. Similarly, while the online channel study (Study 3) is situated in an early-adoption setting, the platform analysis (Study 4) focuses on a single period of pandemic disruption. These contextual boundaries invite replication across other geographies, sectors, and timeframes to test the robustness of the patterns identified here. Future research could extend the framework to new categories or retail systems to examine whether the interaction between context, strategy, and consumer response persists under different market conditions (\cite{Gielens2019-si}).

A second set of limitations concerns the focus and operationalization of key constructs. The pricing studies isolate regular price and discount effects but do not include such other marketing instruments as in-store displays, advertising, and supplier promotions (\cite{Bijmolt2005-kg,Nijs2001-qq}). This simplification allows for clean estimation of price responsiveness but may understate the joint influence from other marketing actions. In the online channel study, the quasi-experimental design captures behavioral outcomes following adoption, but it simplifies the adoption process itself. Channel adoption is modeled as a discrete event by abstracting from managerial interventions (e.g., communication campaigns or delivery incentives) that might influence both the likelihood and intensity of adoption. Moreover, the focus on grocery purchases, with their unique delivery constraints and purchase rhythms, limits the generalizability of behavioral patterns to other product or service categories (\cite{Campo2020-xu,Kushwaha2013-iu}). Finally, the platform study uses playlist followers as a proxy for engagement, an approach that captures visibility and attention but not actual consumption or revenue outcomes. Together, these modeling constraints highlight a trade-off between analytical clarity and strategic completeness. Future research should address this gap by integrating richer managerial, cost, and behavioral data to achieve a more comprehensive understanding of performance outcomes.

A final limitation stems from reliance on observational data. The quasi-experimental design of Study 3 and the fixed effects approach in Study 4 strengthen causal interpretation but cannot fully eliminate unobserved heterogeneity or selection bias (\cite{Goldfarb2022-jt}). Early adopters may differ systematically from non-adopters in ways that are difficult to observe, from digital literacy to risk tolerance and time constraints, which can influence both their likelihood of adoption and their subsequent purchasing behavior (\cite{Avery2012-ex,Li2015-tq,Rietveld2018-ay}). Similarly, playlist performance may reflect algorithmic or editorial choices that, made by the platform, are not visible to researchers and make it challenging to disentangle organic audience behavior from platform-driven amplification. Across all four studies, the reliance on observational data brings the typical trade-offs of real-world research: high external validity but limited causal precision. Future research could combine field experiments, A/B testing, and simulation-based methods with observational data to more directly examine the underlying behavioral processes that cannot be observed in the current data. Integrating survey-based or qualitative insights could also reveal psychological mechanisms behind the behavioral patterns observed.

Overall, the limitations identified in this thesis highlight opportunities for future research. Addressing these through broader contexts, richer data, and complementary methods would help strengthen and extend the insights developed here.
